\documentclass[11pt]{article}
\usepackage[margin=1in]{geometry}
\usepackage{hyperref}
\usepackage{xcolor}
\usepackage{longtable}
\usepackage{array}
\usepackage{enumitem}
\definecolor{goalosblue}{HTML}{0B1020}
\definecolor{goalosgold}{HTML}{B78A00}
\hypersetup{colorlinks=true, linkcolor=goalosblue, urlcolor=goalosblue}
\title{\textbf{AEP-008 --- Proof Room Standard}\\\large Governed Workspaces for Institutional Machine-Work Proof}
\author{Vincent Boucher\\QUEBEC.AI \& MONTREAL.AI}
\date{Version 1.1 Institutional Edition --- 2026-06-05}
\begin{document}
\maketitle

\begin{abstract}
AEP-008 defines the Proof Room: the governed workspace where machine work becomes institutional proof. A Proof Room binds commitments, runs, ProofPackets, Evidence Dockets, Selection Gate decisions, Tool Permission records, Rollback Receipts, and Public-Safe Proof Reports into one controlled room with roles, scope, boundaries, review, approvals, publication rules, audit export, and closure.
\end{abstract}

\section*{Canonical Thesis}
A model can answer. An agent can act. An institution must prove. A network must select. A tool must be permitted. A release must recover. A public claim must be safe. A Proof Room makes the system operable.

\section*{Canonical Law}
No room without charter.\\
No run without scope.\\
No evidence without boundary.\\
No promotion without gate.\\
No tool without permission.\\
No release without rollback.\\
No public proof without redaction.\\
No closure without archive.\\
No proof, no evolution.

\section{Proof Room Modes}
\begin{longtable}{p{0.23\linewidth}p{0.67\linewidth}}
\textbf{Mode} & \textbf{Purpose}\\
sandbox & Low-risk learning, demonstration, internal experimentation, and educational use.\\
team & Shared team workflow improvement and team-level proof.\\
institutional & Enterprise, regulated, audit-ready, or internal governance use.\\
public\_safe & Preparation and approval of external proof reports.\\
incident & Recovery, rollback, quarantine, compensation, and post-incident review.\\
sovereign & Public-sector, national, regulated, jurisdictional, or protected institutional use.\\
\end{longtable}

\section{Required Objects}
AEP-008 defines Proof Room Manifest, Proof Room Charter, Scope Boundary, Role Assignment Registry, Evidence Boundary, Work Item Registry, Room Session Record, Artifact Registry, Tool Permission Register, Evidence Docket Register, ProofPacket Register, Selection Gate Register, Rollback Register, Public Report Register, Decision Log, Room Audit Export, and Room Closure Report.

\section{Conformance Levels}
\begin{longtable}{p{0.20\linewidth}p{0.70\linewidth}}
Level 0 & Informal Proof Room.\\
Level 1 & Basic Proof Room with manifest, charter, scope, roles, evidence boundary, and work item registry.\\
Level 2 & Operational Proof Room with docket, packet, tool permission, and decision registers.\\
Level 3 & Institutional Proof Room with gates, rollback, public reports, separation of duties, audit export, and closure report.\\
Level 4 & Regulated Proof Room with protected boundary, retention, jurisdiction, approvers, independent review, and audit export.\\
Level 5 & Sovereign Proof Room with sovereign jurisdiction, public/private/protected controls, public-institution roles, policy traceability, reporting controls, and archive requirements.\\
\end{longtable}

\section{Claim Boundary}
AEP-008 does not claim achieved AGI, achieved ASI, perfect safety, legal compliance certification, financial or legal advice, guaranteed ROI, production readiness, government endorsement, or national-security readiness. It defines an operating standard for Proof Rooms.

\section*{Canonical Public Line}
AEP-001 defines the protocol. AEP-002 defines the docket. AEP-003 defines the packet. AEP-004 defines the gate. AEP-005 defines the permission. AEP-006 defines recovery proof. AEP-007 defines public proof. AEP-008 defines the room. GoalOS turns machine work into governed, reviewable, recoverable, and public-safe institutional proof.

\end{document}
